\chapter{Literature Review}
\label{chap:literature}

This chapter reviews the established methods on which \NEMO\ is based. The
review covers constrained structural optimization, derivative-free search,
Nelder--Mead, variable scaling, Latin-hypercube sampling, tiered model
fidelity, CAD--CAE integration, simulation intent, marine structural
applications, and computational-model credibility. It then relates these
topics to the bracket, padeye, and stabilizer case studies and defines the
limited research gap addressed by this project.

\section{Structural Optimization Methods}
\label{sec:structural-optimization}

\subsection{Constrained Search}

Structural optimization formalizes the selection of design variables to
improve an objective while satisfying constraints. An objective can represent
mass, cost, compliance, displacement, stress, fatigue damage, environmental
impact, or a weighted combination of measures. Constraints can represent
strength, stiffness, stability, vibration, geometry, fabrication, inspection,
or regulatory limits. The problem is written in the general form

\begin{align}
\underset{\mathbf{x}}{\text{minimize}}\quad
    & f(\mathbf{x}), \label{eq:general-objective}\\
\text{subject to}\quad
    & g_j(\mathbf{x})\leq 0,\quad j=1,\ldots,m,\\
    & h_k(\mathbf{x})=0,\quad k=1,\ldots,p,\\
    & \mathbf{l}\leq\mathbf{x}\leq\mathbf{u},
\end{align}

where $\mathbf{x}$ is the design vector, $f$ is the objective, $g_j$ and
$h_k$ are inequality and equality constraints, and $\mathbf{l}$ and
$\mathbf{u}$ are variable bounds.

Structural optimization is commonly divided into size, shape, and topology
optimization. Size optimization changes dimensions such as thickness and
section area. Shape optimization changes geometric boundaries while retaining
a general feature arrangement. Topology optimization changes the distribution
or connectivity of material inside a design region. These classes can overlap,
but the distinction helps define what an optimization claim means.

\NEMO\ performs size and bounded parametric shape optimization. The number and
general arrangement of major features remain fixed. The optimizer changes
dimensions such as plate thickness, rib height, lug width, spar position, and
fillet radius. It does not remove arbitrary material or create a new topology.
This restriction keeps every candidate compatible with an intended load path,
a native CAD generator, and predefined load and support roles.

The project uses mass as the principal objective. This choice is easy to
measure and compare, but it does not make mass the only design criterion. A
low-mass candidate is useful only if its constraints pass and its omitted
failure modes are later checked. In marine practice, this qualification is
especially important because cyclic loading, corrosion, fabrication detail,
and uncertain operational loads can govern a design that passes one static
yield check.

\subsection{Derivative-Free Optimization}
\label{sec:derivative-free}

Many engineering analyses appear as black-box functions to the optimizer. The
optimizer supplies dimensions and receives calculated outputs, but it does not
have exact derivatives of those outputs. CAD regeneration, meshing changes,
solver tolerances, discrete feature behavior, and failed candidates can also
make numerical gradients unreliable.

Direct-search and derivative-free methods use objective values without
requiring analytical gradients. Kolda, Lewis, and Torczon reviewed the
development of direct-search methods and distinguished methods with formal
convergence properties from more heuristic searches \cite{kolda2003}. Their
review shows that derivative-free operation is not the same as assumption-free
operation. The designer must still scale variables, handle bounds, define
constraints, select termination criteria, and evaluate sensitivity to the
starting point.

Derivative-free methods can be attractive when each function evaluation is
moderately expensive and the number of variables is limited. They are less
attractive when the dimension is very high or when gradients can be calculated
cheaply and reliably. The three \NEMO\ cases contain six or nine continuous
variables. Their analytical evaluations are inexpensive, while the wider
workflow includes potential discontinuities from geometry checks. This
combination supports the use of a simple derivative-free method for the
reference implementation.

\subsection{The Nelder--Mead Simplex Method}
\label{sec:nelder-mead-review}

Nelder and Mead introduced the simplex method in 1965 \cite{nelder1965}. In an
$n$-dimensional space, the method maintains $n+1$ vertices. It evaluates and
ranks these vertices, calculates the centroid of the better vertices, and
attempts to replace the worst vertex. The standard trial operations are
reflection, expansion, contraction, and shrink.

If $\mathbf{x}_h$ is the worst vertex and $\bar{\mathbf{x}}$ is the centroid
of the remaining vertices, a reflected point is

\begin{equation}
\mathbf{x}_r=\bar{\mathbf{x}}+
\alpha(\bar{\mathbf{x}}-\mathbf{x}_h),
\label{eq:nm-reflection}
\end{equation}

where $\alpha$ is the reflection coefficient. A sufficiently good reflection
can lead to expansion:

\begin{equation}
\mathbf{x}_e=\bar{\mathbf{x}}+
\gamma(\mathbf{x}_r-\bar{\mathbf{x}}).
\label{eq:nm-expansion}
\end{equation}

An unsuccessful reflection can cause contraction toward the centroid. When
contraction also fails, the simplex shrinks toward its best vertex. The
standard coefficients used by \NEMO\ are $\alpha=1$, $\gamma=2$, $\rho=0.5$,
and $\sigma=0.5$.

The method has several practical advantages. It is compact, does not require a
gradient, and is easy to explain and implement. It also has important limits.
Lagarias \emph{et al.} established convergence results under restricted
conditions and discussed cases where the method can approach a non-minimizer
\cite{lagarias1998}. Nelder--Mead is a local search. Its final design can
depend on the initial simplex, variable scaling, bounds, penalties, and the
shape of the objective landscape.

These limitations affect the language of this report. One local run cannot
support a claim of global optimality. \NEMO\ uses a seeded design-space sample
and three starts to reduce dependence on the configured baseline. The report
uses the phrase ``best feasible design found within this parameterized
search.'' This wording is more defensible than ``global optimum.''

\begin{table}[H]
\centering
\caption{Nelder--Mead operations and their purpose in the implemented search}
\label{tab:nm-operations}
\begin{tabularx}{\textwidth}{p{0.18\textwidth}p{0.18\textwidth}X}
\toprule
\textbf{Operation} & \textbf{Coefficient} & \textbf{Purpose}\\
\midrule
Reflection & $\alpha=1$ & Tests a point across the centroid from the current
worst vertex.\\
Expansion & $\gamma=2$ & Extends a successful reflected step when it improves
the best objective.\\
Contraction & $\rho=0.5$ & Tests a shorter step when reflection does not give
adequate improvement.\\
Shrink & $\sigma=0.5$ & Moves all non-best vertices toward the best point after
failed trial steps.\\
\bottomrule
\end{tabularx}
\end{table}

\subsection{Variable Scaling, Bounds, and Constraint Penalties}
\label{sec:scaling-review}

Engineering variables can differ greatly in magnitude and unit. A plate
thickness may be 4~mm, a chordwise spar position may be 30\%, and a root
insert length may be 180~mm. If an optimizer works directly in these physical
coordinates, its simplex geometry and trial step sizes can be dominated by
the numerically largest variables.

\NEMO\ maps each physical variable $x_i$ to the scaled coordinate

\begin{equation}
z_i=\frac{x_i-l_i}{u_i-l_i},
\qquad 0\leq z_i\leq 1,
\label{eq:variable-scaling}
\end{equation}

where $l_i$ and $u_i$ are the lower and upper bounds. The evaluator performs
the inverse mapping before it uses the engineering equations. Trial points are
clipped at the scaled bounds.

Scaling gives each coordinate the same numerical range. It does not make the
physics equally sensitive. A small thickness change can alter both mass and
section properties strongly, while a similar scaled change in a fillet radius
may affect only one stress factor in the analytical model. Sensitivity remains
a property of the model.

Nelder--Mead does not directly enforce general nonlinear constraints in this
implementation. \NEMO\ therefore uses a penalty objective. If
$m(\mathbf{x})$ is mass, $F(\mathbf{x})$ is factor of safety,
$F_{\min}$ is the minimum factor of safety, $\delta(\mathbf{x})$ is
displacement, and $\delta_{\max}$ is the displacement limit, the objective is

\begin{equation}
\begin{split}
J(\mathbf{x})=m(\mathbf{x})+w\Bigg[
&\left(
\frac{\max(0,F_{\min}-F(\mathbf{x}))}{F_{\min}}
\right)^2\\
&+\left(
\frac{\max(0,\delta(\mathbf{x})-\delta_{\max})}
{\delta_{\max}}
\right)^2
\Bigg],
\end{split}
\label{eq:penalty-objective}
\end{equation}

with $w=100$ for the current parts.

A feasible design has $J=m$. An infeasible design has a higher objective
according to its normalized violations. Penalties simplify the search, but
their weight affects ranking. A weak penalty can make an infeasible light
design more attractive than a feasible heavier design. A large penalty can
create a steep landscape near the boundary. The final selection must therefore
filter explicitly for feasibility. The stabilizer case in
Chapter~\ref{chap:results} demonstrates why this check is necessary.

\subsection{Latin-Hypercube Sampling}
\label{sec:lhs-review}

McKay, Beckman, and Conover introduced Latin-hypercube sampling as a stratified
method for computer experiments \cite{mckay1979}. For each variable, an
$N$-point Latin hypercube samples once from each of $N$ equal-probability
intervals. Independent permutations combine the one-dimensional samples into
$N$ multi-dimensional points.

Compared with an unstructured random sample of the same size, a Latin
hypercube improves one-dimensional stratification. It does not guarantee dense
coverage of every interaction, especially in high dimensions. Sixty points
remain sparse in a nine-dimensional space. The method is therefore best
viewed here as an economical exploration step.

\NEMO\ uses 60 samples per part with seed 42. The fixed seed makes the
campaign repeatable. Sampling serves three purposes:

\begin{enumerate}
    \item it checks whether calculated outputs change in a physically sensible
    direction over the bounds;
    \item it identifies low-mass feasible candidates that can start local
    search; and
    \item it reduces reliance on one baseline-start Nelder--Mead run.
\end{enumerate}

Sampling also provides diagnostic evidence. A high failure rate can reveal
poor parameter bounds or a fragile geometry generator. An unexpectedly high
feasibility rate can reveal an over-conservative baseline or a screening model
that does not activate the declared constraints.

\section{Model Fidelity and CAD--CAE Integration}
\label{sec:model-fidelity-integration}

\subsection{Tiered Model Fidelity}
\label{sec:fidelity-review}

Engineering optimization often combines models with different cost and
accuracy. Peherstorfer, Willcox, and Gunzburger review multi-fidelity methods
that combine inexpensive low-fidelity models with expensive high-fidelity
models \cite{peherstorfer2018}. Formal multi-fidelity optimization uses the
high-fidelity model inside the optimization strategy to improve prediction,
selection, or convergence.

\NEMO\ does not yet implement such a closed loop. The analytical model guides
the search. Native CAD verifies geometry and mass for selected designs. Manual
FEA follows as a separate validation stage. The correct description is a
\emph{tiered-fidelity workflow}. The analytical optimum is not automatically
corrected by FEA.

\begin{figure}[H]
\centering
\begin{tikzpicture}[
    node distance=0.75cm,
    every node/.style={font=\small},
    stage/.style={
        rectangle,
        rounded corners,
        draw=nemoblue,
        fill=nemolight,
        thick,
        minimum width=0.83\textwidth,
        minimum height=1.25cm,
        align=center
    }
]
\node[stage] (a) {\textbf{Tier 1: Analytical screening}\\
Low cost; many candidates; first-order mass, stress, FOS, and displacement};
\node[stage,below=of a] (b) {\textbf{Tier 2: Native CAD verification}\\
Selected candidates; solid validity, physical mass, exports, and boundary roles};
\node[stage,below=of b] (c) {\textbf{Tier 3: Manual FEA verification}\\
Few candidates; checked material, loads, supports, mesh, reactions, and results};
\node[stage,below=of c] (d) {\textbf{Tier 4: Design qualification}\\
Codes, detailed failure modes, uncertainty, physical tests, and engineering approval};
\draw[arrow] (a) -- (b);
\draw[arrow] (b) -- (c);
\draw[arrow] (c) -- (d);
\end{tikzpicture}
\caption{Tiered evidence model used by the project}
\label{fig:tiered-evidence}
\end{figure}

The distinction between the tiers protects the report from a common error:
treating agreement on one quantity as complete validation. Close agreement
between analytical and CAD mass shows that the volume approximations represent
the generated geometry reasonably well. It does not show that the analytical
stress model represents local three-dimensional stress.

\subsection{CAD--CAE Integration}
\label{sec:cad-cae-review}

Automated structural optimization requires dependable communication between
geometry, analysis, and optimization. Park and Dang coupled commercial CAD and
CAE software through scripting, programming interfaces, and metamodels
\cite{park2010}. Their work shows how integration can reduce repetitive work
and analysis cost. It also introduces an additional source of error when a
surrogate model approximates the high-fidelity response.

Li \emph{et al.} developed a unified parametric representation based on
extended voxels. Their method supports a closed loop among design, simulation,
and optimization while carrying design and simulation meaning
\cite{li2023}. This work confirms that automated CAD--CAE optimization is an
established research area. \NEMO\ does not claim novelty from automation
alone.

The project instead uses a modest file-based architecture. External Python
handles engineering definitions, optimization, run records, and validation
packages. Fusion handles native CAD and physical properties. Versioned JSON
files connect the two environments. This approach is slower than an in-process
call but has useful properties:

\begin{itemize}
    \item a request and response can be inspected with an ordinary text editor;
    \item the Fusion environment does not need third-party optimization
    packages;
    \item a failed request can preserve a controlled error record;
    \item the protocol can retain compatibility with the original bracket
    proof of concept; and
    \item the same external application can later support another CAD or FEA
    backend.
\end{itemize}

The architecture also has limits. A single shared request directory serializes
CAD work. Polling adds latency. Local file locking requires retry behavior.
The bridge cannot make an unavailable solver API available. These are accepted
trade-offs for an inspectable student-scale implementation.

\subsection{Simulation Intent and Semantic Boundary Roles}
\label{sec:simulation-intent}

Geometry changes can invalidate an analysis even when the regenerated solid
looks correct. A load attached to a low-level face identifier can disappear or
move after a feature changes the boundary representation. A reliable
parametric workflow must preserve why a region exists and how it is used.

Nolan \emph{et al.} define simulation intent as the high-level modeling and
idealization decisions required for fit-for-purpose analysis
\cite{nolan2015}. Their work associates analysis meaning with regions and
interfaces rather than relying only on unstable topology.

\NEMO\ applies the same principle in a smaller implementation. The bracket
declares \code{fixed\_support} and \code{equipment\_load}. The padeye declares
\code{fixed\_support} and \code{pin\_bearing}. The stabilizer declares
\code{fixed\_support}, four pressure bands, and \code{tip\_monitor}. The
Fusion generator selects faces with part-specific geometric rules and records
area, centroid, normal, bounding box, and cylinder data where available.

Semantic roles improve traceability, but they do not remove the need for human
inspection. A broad geometric selector can match an unintended face. A
topology change can divide one intended region into several faces. The manual
FEA procedure must display and confirm every selected face before solving.

\section{Marine Structural Applications}
\label{sec:marine-optimization-review}

Marine structural optimization must represent realistic loads, supports, and
failure modes. Gentils, Wang, and Kolios optimized an offshore wind support
structure with parametric FEA and a genetic algorithm
\cite{gentils2017}. Their constraints included stress, displacement,
vibration, buckling, and fatigue. The study reported a 19.8\% mass reduction
for the complete support structure.

Zheng, Li, and Xiao later used sensitivity analysis, surrogate models, and a
genetic algorithm for jacket structures \cite{zheng2023}. Their method
addressed natural frequency, displacement, stress, buckling, and later fatigue
checks. Wang, Mantey, and Zhang developed a fast parametric FEA and
genetic-algorithm tool for jacket structures \cite{wang2024}. These studies
show that marine structural optimization commonly requires several load cases
and constraints beyond static yield.

The literature provides two lessons for \NEMO. First, mass reduction must
remain subordinate to the governing engineering constraints. Second, the
selected design can change when supports and loads change. The project
therefore treats its load cases as declared engineering envelopes, not as
complete vessel or offshore design cases.

\begin{table}[H]
\centering
\caption{Relationship between selected literature and the NEMO implementation}
\label{tab:literature-comparison}
\begin{tabularx}{\textwidth}{p{0.23\textwidth}XX}
\toprule
\textbf{Source area} & \textbf{Established lesson} & \textbf{Use in NEMO}\\
\midrule
Direct search \cite{kolda2003,nelder1965,lagarias1998} &
Derivative-free local search is useful but needs scaling, stopping rules, and
careful claims. &
Scaled bounded variables, three starts, and ``best design found'' wording.\\

Computer experiments \cite{mckay1979} &
Stratified samples improve economical exploration. &
Sixty seeded Latin-hypercube samples per part.\\

Multi-fidelity methods \cite{peherstorfer2018} &
Models of different cost require a defined relationship. &
Tiered analytical, CAD, manual FEA, and qualification evidence.\\

CAD--CAE integration \cite{park2010,li2023} &
Automation must preserve data and analysis meaning. &
Versioned JSON bridge, native CAD, and explicit partial responses.\\

Simulation intent \cite{nolan2015} &
Boundary meaning must survive geometry change. &
Semantic roles and geometric face signatures.\\

Marine optimization \cite{gentils2017,zheng2023,wang2024} &
Mass is one objective among strength, stiffness, stability, dynamics, and
fatigue. &
Static screening followed by a stated list of required detailed checks.\\
\bottomrule
\end{tabularx}
\end{table}

\subsection{Bracket Structural Behavior}
\label{sec:bracket-review}

Autodesk lists brackets among components suited to linear Static Stress
analysis when the assumptions of small deformation, unchanged load direction,
and linear elastic material behavior hold \cite{autodeskStaticStress}. The
apparent simplicity of a bracket can be misleading. Local stress depends on
load introduction, fastener restraint, contact, rib geometry, fillet radius,
plate flexibility, and weld detail.

The \NEMO\ bracket is an equipment-foundation archetype with a baseplate, four
mounting holes, two ribs, and rib-root fillets. Its analytical model uses an
equivalent cantilever relation. This model captures the broad effects of
section depth, plate thickness, rib thickness, and load eccentricity. It does
not resolve bolt bearing, fastener preload, prying, weld-toe stress, local
contact, or three-dimensional restraint.

The bracket is useful as a proof of concept because the geometry is easy to
recognize and the load path is easy to explain. It is not useful because it is
free from modeling risk. The proof-of-concept FEA report later showed that a
three-order-of-magnitude input error can survive through a completed solver
run. This finding supports a workflow that checks input values and reaction
balance before it discusses result contours.

\subsection{Padeyes and Pinned Connections}
\label{sec:padeye-review}

Padeyes transfer lifting force through a pin hole into a lug, reinforcement,
welds, and supporting structure. Potential governing effects include
pin-bearing stress, net-section tension, shear-out, bending, local contact,
lug-root stress, gusset load transfer, and weld stress. Clearance, load angle,
out-of-plane load, and local geometry can strongly affect the stress field.

Soh and Soh compared simplified two-dimensional padeye models with a
three-dimensional finite-element solution \cite{soh1989}. Their work showed
that model selection affects stress prediction and can lead to conservative
overdesign. Liu, Zhou, and Tan used FEA to study padeye stress and deformation
and then formulated an optimization problem to reduce steel use
\cite{liu2013}. These studies support the use of a simplified model for early
screening, provided later FEA represents the important three-dimensional
behavior.

Pedersen examined the stress concentration around a pin-loaded hole and
showed the influence of contact modeling and local geometric refinement
\cite{pedersen2019}. This result is directly relevant to the \NEMO\ padeye.
Its current analytical model compares several nominal stress modes, but it
does not model a separate pin, clearance, nonlinear contact, or a weld group.
A later FEA campaign must decide whether a distributed bearing load is
adequate for screening or whether explicit pin contact is required.

Marine lifting appliances also require a governing rule set. DNV-ST-N001
addresses marine operations and associated structural considerations
\cite{dnvN001}. Lloyd's Register publishes a Code for Lifting Appliances in a
Marine Environment that covers design, materials, fabrication, testing,
marking, survey, and documentation \cite{lrLifting2026}. \NEMO\ does not
claim compliance with either source. They are cited to show the wider
qualification work that lies beyond a single yield-based optimization.

\subsection{Stabilizer Fins and Prescribed Hydrodynamic Loads}
\label{sec:stabilizer-review}

The stabilizer differs from the bracket and padeye because its load is
distributed over an external lifting surface. The project retains a NACA 0015
section. Abbott, von Doenhoff, and Stivers documented NACA airfoil families
and aerodynamic data in NACA Report 824 \cite{abbott1945}. The use of a known
profile makes the geometric definition reproducible, but it does not supply a
complete marine pressure field.

The \NEMO\ stabilizer keeps its span, root chord, tip chord, sweep, trailing
edge, root flange, and rib stations fixed. It changes skin thickness, two spar
positions, two spar thicknesses, rib thickness, root-insert dimensions, and
root fillet radius. Its prescribed resultant is calculated from a simplified
lift expression and distributed through four spanwise weights.

This setup supports internal structural sizing under one declared load
envelope. It does not optimize hydrodynamic efficiency. It omits vessel
motions, unsteady flow, free-surface effects, ventilation, cavitation,
actuator behavior, pressure reversal, impact, and fluid--structure
interaction. A later design stage must derive the pressure from an accepted
hydrodynamic method, CFD, model testing, or full-scale data.

\section{Finite Element Verification and Validation}
\label{sec:linear-static-fea-review}

\subsection{Linear Static Analysis}

Autodesk describes a Static Stress study as appropriate when deformation is
small, load direction does not change significantly, material properties are
constant, and the response remains linear elastic
\cite{autodeskStaticStress}. A complete setup requires geometry, materials,
constraints, loads, contacts, mesh, solution, and result review
\cite{autodeskStaticSetup}.

These requirements are coupled. A correct material with a wrong load gives a
wrong result. A correct load on the wrong face can change the load path. An
overly rigid support can reduce deformation and create a local stress
singularity. A coarse mesh can underpredict a local stress. A maximum contour
value can occur at a mathematical singularity that is unsuitable for direct
acceptance.

Reaction forces provide a basic equilibrium check. Autodesk notes that support
reactions should oppose the total applied load \cite{autodeskReaction}. For
the bracket, a correctly configured 1472~N downward load should produce a
vertical support reaction close to 1472~N, subject to sign convention and
small numerical differences. A reaction near 1.472~N would immediately expose
the factor-of-one-thousand load error.

Mesh convergence is also necessary. Autodesk explains that coarse meshes
commonly underpredict stress and that stress is more mesh-sensitive than
displacement \cite{autodeskMesh}. Comparison should use the same geometric
location across refinements. A credible record should report the global mesh,
local refinements, element order, node and element count, the monitored
quantity, and the change between successive meshes.

\subsection{Verification, Validation, and Model Credibility}
\label{sec:vv-review}

Verification asks whether equations and software were solved or implemented
correctly. Validation asks whether the model represents the real system well
enough for its intended use. These questions are related but distinct. A
correctly implemented simplified model can still be unsuitable for a service
decision.

ASME V\&V 10 provides a framework and terminology for verification,
validation, and uncertainty quantification in computational solid mechanics
\cite{asmeVV10}. NASA-STD-7009B emphasizes intended use, acceptance criteria,
input pedigree, uncertainty, sensitivity, verification, validation, and
communicated credibility \cite{nasa7009}. The principles are relevant to
\NEMO\ even though the project does not fall under NASA governance. Model
credibility depends on the decision being supported.

\begin{table}[H]
\centering
\caption{Evidence levels and the questions they answer}
\label{tab:evidence-questions}
\begin{tabularx}{\textwidth}{p{0.20\textwidth}XX}
\toprule
\textbf{Evidence level} & \textbf{Primary question} & \textbf{Representative
checks}\\
\midrule
Software verification &
Did the code implement its declared contracts and calculations consistently? &
Unit tests, bounds, baseline feasibility, failure handling, logging schema,
repeatable seeds.\\

CAD verification &
Did the requested dimensions create a usable native solid and required
artifacts? &
One connected solid, positive volume, STEP export, boundary roles, physical
mass.\\

FEA verification &
Was the numerical study configured and solved correctly? &
Material, load, support, contact, reaction balance, mesh convergence, result
location.\\

Physical validation &
Does the model represent measured component behavior adequately for its
intended use? &
Proof load, strain, displacement, failure mode, repeatability, uncertainty.\\

Design qualification &
Does the complete design satisfy governing service requirements? &
Design codes, load combinations, fatigue, buckling, corrosion, fabrication,
inspection, competent approval.\\
\bottomrule
\end{tabularx}
\end{table}

The table shows why a project can complete software and CAD verification
without completing structural validation. It also explains why a single FEA
contour is not final proof. The bracket report is retained because it documents
the setup and exposes a process weakness. Its incorrect numerical result is
not promoted to a higher evidence level.

\section{Literature Synthesis and Research Gap}
\label{sec:literature-synthesis}

The literature establishes every main method used by \NEMO. Nelder--Mead
provides a compact derivative-free local search
\cite{nelder1965,lagarias1998}. Direct-search literature explains the need for
careful implementation and claims \cite{kolda2003}. Latin-hypercube sampling
provides a repeatable exploration stage \cite{mckay1979}. Multi-fidelity
literature explains the value and limits of combining models of different
cost \cite{peherstorfer2018}. CAD--CAE research identifies integration and
simulation intent as central problems \cite{park2010,nolan2015,li2023}.
Marine optimization studies show that useful lightweighting depends on
several constraints and credible load models
\cite{gentils2017,zheng2023,wang2024}. Computational-mechanics guidance
requires a declared intended use and an evidence chain
\cite{asmeVV10,nasa7009}.

No single method in \NEMO\ is new. The reviewed literature also contains more
advanced automated loops, global optimizers, surrogate models, and
high-fidelity marine applications. The project therefore does not claim to be
the first marine optimizer or the first CAD--CAE automation system.

The limited gap addressed here is practical integration at student-project
scale. The reviewed sources do not directly provide a transparent,
part-manifest-driven workflow that combines unit-aware definitions,
first-order screening, seeded sampling, scaled bounded multi-start search,
failure-contained evaluation, versioned atomic file exchange, native Fusion
CAD, semantic boundary metadata, validation packages, and explicit separation
of analytical, CAD, and FEA evidence across a bracket, padeye, and stabilizer.

\NEMO\ addresses that implementation gap as an auditable reference workflow.
Its value lies in exposing each assumption and handoff, reducing repetitive
candidate work, and preserving enough evidence for a supervisor or engineer to
check the result. The next chapter describes how the workflow was implemented
and evaluated.
